\documentclass[aspectratio=169]{beamer}

\usetheme{Madrid}
\usecolortheme{default}

\usepackage[T1]{fontenc}
\usepackage[utf8]{inputenc}
\usepackage[french]{babel}
\usepackage{lmodern}
\usepackage{booktabs}
\usepackage{hyperref}

	itle[JPM Native]{JPM --- Moteur natif (Phase 2)}
\subtitle{Simulation Scrum (Gestion de Projet) : blocages, correctifs, performances vs Maven}
\author{Équipe projet (go-jpm)}
\date{26 décembre 2025}

\begin{document}

\begin{frame}
  \titlepage
\end{frame}

\begin{frame}{Contexte (Gestion de Projet) \& objectifs}
\begin{itemize}
  \item Simulation Scrum : enseignant = manager/validateur, équipe = dev
  \item Produit : CLI \texttt{jpm} (Go) pour projets Java
  \item Objectif technique : \textbf{moteur natif} (sans \texttt{mvn})
  \item Objectif projet : incréments livrables + traçabilité (tests, docs, bench)
\end{itemize}
\end{frame}

\begin{frame}{Ce que nous avons construit}
\begin{itemize}
  \item Pipeline natif : \texttt{resolve} $\rightarrow$ \texttt{fetch} $\rightarrow$ \texttt{compile} $\rightarrow$ \texttt{package}
  \item Résolution Maven : POM + transitifs (graphe)
  \item Téléchargement Maven Central + cache disque
  \item Compilation via \texttt{javac}
  \item Packaging JAR déterministe + lockfile \texttt{jpm.lock.yaml}
\end{itemize}
\end{frame}

\begin{frame}{Principaux blocages (POMs réels) et correctifs}
\begin{enumerate}
  \item \textbf{Interpolation de propriétés}: versions \texttt{${hamcrestVersion}}
    \begin{itemize}
      \item Correctif : parse \texttt{<properties>} + substitution \texttt{${...}}
    \end{itemize}
  \item \textbf{Encodages non UTF-8}: ISO-8859-1 déclaré
    \begin{itemize}
      \item Correctif : XML decoder + conversion d’encodage
    \end{itemize}
  \item \textbf{Versions manquantes}: dépendances gérées par \texttt{dependencyManagement}
    \begin{itemize}
      \item Symptôme : coordonnées invalides (version vide) $\rightarrow$ retries HTTP
      \item Correctif : filtrer les deps sans version (en attendant DM)
    \end{itemize}
\end{enumerate}
\end{frame}

\begin{frame}{Performance : protocole de benchmark}
\begin{itemize}
  \item Fixtures : \texttt{bench/projects} (small/medium/large)
  \item \texttt{RUNS=10} ; métriques \textbf{warm} = runs 2..10
  \item Comparaison Maven : ``bridge'' (\texttt{mvn package}) sur copie temporaire
\end{itemize}
\end{frame}

\begin{frame}{Performance : résultats (RUNS=10, 20:42 UTC)}
\small
\begin{tabular}{lrrr}
\toprule
	extbf{Projet} & \textbf{Native warm mean (ms)} & \textbf{Maven warm mean (ms)} & \textbf{Gagnant}\\
\midrule
small  & 391  & 2021 & Native\\
medium & 617  & 2374 & Native\\
large  & 5307 & 2692 & Maven\\
\bottomrule
\end{tabular}

\vspace{0.5em}
\normalsize
\begin{itemize}
  \item Gain majeur : suppression des retries sur coordonnées invalides
  \item Ajout : téléchargement d’artefacts en parallèle (pool borné)
\end{itemize}
\end{frame}

\begin{frame}{Pourquoi Maven gagne encore sur \textit{large}}
\begin{itemize}
  \item Le coût restant est dominé par la compilation à l’échelle
  \item Classpath peut contenir plusieurs versions d’un même artefact
  \item Maven : compilation incrémentale + gestion de versions mature
\end{itemize}
\end{frame}

\begin{frame}{Backlog (si plus de temps)}
\begin{itemize}
  \item Support \texttt{dependencyManagement}
  \item Déduplication/choix de version effectif du classpath (nearest-wins)
  \item Paralléliser la résolution POM (concurrence bornée)
  \item Compilation incrémentale
\end{itemize}
\end{frame}

\begin{frame}{Reproductibilité}
\begin{itemize}
  \item Bench: \texttt{RUNS=10 bench/bench.sh}
  \item Résultats: \texttt{bench/results/latest/bench.json}
\end{itemize}
\end{frame}

\begin{frame}
  \centering
  \large Questions?
\end{frame}

\end{document}
